% ------------------------------------------------------------------
%  Results (bold run-in subheadings that state
%  the finding, present tense for figure description, past tense for
%  what was done, no methods detail).
%
%  NUMBERS: LongPhase 2 values (SNV-only and four-class, with and
%  without GNN; phased fraction, switch errors, Hamming, N50, indel
%  fraction) were replaced on 2026-09-20 with the exact values from the
%  `longphase compare` tables in "GNN source/sw_*.log" (v5.0q, chr1-22;
%  10x has nine replicates there, seed 10 missing). WhatsHap, HapCUT2,
%  LongPhase 1.0, the SNV+indel/SV/5mC configurations, SV and 5mC phased
%  fractions, runtimes and the v4.2.1 counts are still read from the
%  plots in LongPhaseGNN_0902_2-1.pptx (slides 7-11).
%
%  REPLICATE STRUCTURE: 10--20x = mean of ten down-sampling replicates;
%  30--60x = a single replicate. Every quoted value must carry the
%  corresponding dispersion (10--20x) or a robustness statement (30--60x).
% ------------------------------------------------------------------

\subsection*{\toolname co-phases SNVs, indels, SVs and methylation in one graph and scores every phased variant}
\toolname resolves haplotypes in three stages (Fig.~\ref{fig:overview}). All heterozygous evidence carried by the same long reads, comprising SNVs, small indels, SVs and allele-specific 5mC sites called by the built-in \texttt{modcall} module, enters a single weighted phasing graph per chromosome. Within that graph each observation is calibrated by base quality and by the reliability of its sequence context, so that homopolymers, tandem repeats and copy-number-altered regions contribute less (Fig.~\ref{fig:overview}a,b). A graph neural network then inspects a local window around each phased variant (Fig.~\ref{fig:overview}c). Variants predicted to be misphased are unphased instead of flipped, and any phase set that becomes disconnected is split, so genotypes are preserved and no new phase assignment is ever introduced. The outputs are phased VCFs for every variant class and haplotagged alignments (Fig.~\ref{fig:overview}d).

We evaluated \toolname on HG002 nanopore R10.4.1 whole-genome data down-sampled from 10$\times$ to 60$\times$ coverage (Methods). Throughout, values quoted at 10--20$\times$ are means over ten independent down-sampling replicates, and values quoted at 30--60$\times$ come from a single replicate. Phasing accuracy was scored against two independent truth sets: the GRCh38-based GIAB v4.2.1 benchmark, and the benchmark derived from the complete diploid telomere-to-telomere (T2T) assembly of the same individual, T2T-HG002 (referred to here as v5.0q), which extends evaluation into the segmental duplications, satellite arrays and other difficult regions that v4.2.1 excludes. Unless stated otherwise, results refer to the v5.0q benchmark and to SNVs, the only variant class that all compared tools phase.

\subsection*{\toolname phases SNVs three- to fivefold more accurately and six- to eightfold faster than \whatshap and \hapcut}
On identical Clair3 variant calls and SNV-only phasing, \toolname made several-fold fewer switch errors than \whatshap~2.8 and \hapcut~1.3.4 at every coverage (Fig.~\ref{fig:snv}). Without the neural network, the switch error rate of \toolname was 0.046\% at 30$\times$ and 0.034\% at 60$\times$, against approximately 0.115\% and 0.085\% for \whatshap and 0.15\% and 0.11\% for \hapcut, a 2.5-fold reduction relative to \whatshap and a 3.2-fold reduction relative to \hapcut (Fig.~\ref{fig:snv}b)\todo{\whatshap and \hapcut rates are still read from the plots (the \toolname values are now exact, from the \texttt{compare} tables); the two fold-changes are averages over 30--60$\times$ (\whatshap 2.6$\times$ at 30$\times$ and 2.4$\times$ at 60$\times$; \hapcut 3.3$\times$ and 3.1$\times$); decide whether to quote ranges instead, as elsewhere in this section}. Correction reduced the rate further, to 0.031\% at 30$\times$ and 0.023\% at 60$\times$, so that \toolname with correction made 3.6--3.9-fold fewer switch errors than \whatshap and 4.7--5.0-fold fewer than \hapcut across 30--60$\times$, and 2.6-fold and 3.1-fold fewer at 10$\times$\todo{quote the \whatshap and \hapcut switch error rates at 10$\times$ (approximately 0.225\% and 0.27\%) so that the 2.6- and 3.1-fold claims are verifiable from the reported numbers}. In absolute terms correction helped most where evidence is sparsest: at 10$\times$ the rate fell by 0.039 percentage points, from 0.121\% to 0.082\%, against 0.015 percentage points at 30$\times$ and 0.011 at 60$\times$; in relative terms the reduction was close to one-third at every coverage (31--34\%). The block-wise Hamming distance, which counts misassigned variants rather than switch points and therefore penalizes long-range flips, showed the same ordering: at 30$\times$ the distance was approximately 2.6\% for \whatshap and \hapcut, 2.4\% for \toolname and 1.8\% after correction, and at 10$\times$ the distance fell from 7.7--7.8\% for the two competing tools to 5.3\% after correction (Fig.~\ref{fig:snv}c).

This gain in accuracy did not cost completeness or haplotype length. All tools phased a similar fraction of benchmark heterozygous SNVs, rising from roughly 78\% at 10$\times$ to a plateau above 91\% from 20$\times$ onwards, with \whatshap phasing about one percentage point more than \toolname at every coverage (approximately 92.5\% versus 91.4--91.6\% at 30--60$\times$), consistent with the more conservative read-based correction of \toolname (Fig.~\ref{fig:snv}a). Block N50 grew almost linearly with coverage for all tools and was highest for \toolname (approximately 2.9~Mb at 60$\times$, versus 2.75~Mb for \whatshap and \hapcut); the neural network shortened blocks by 5--9\% at 10--20$\times$ (0.84~Mb versus 0.92~Mb at 10$\times$) and by less than 3\% from 30$\times$ onwards (2.86 versus 2.91~Mb at 60$\times$) (Fig.~\ref{fig:snv}d). \toolname was also substantially faster: whole-genome phasing took 3~min at 10$\times$ and 15~min at 60$\times$, compared with 26 and 95~min for \hapcut and 26 and 115~min for \whatshap, a 6--8-fold speed-up at 60$\times$ (5--8-fold across 30--60$\times$); the neural network added 1--2~min (Fig.~\ref{fig:snv}e)\todo{threads and hardware; exact times, including the 30$\times$ times (approximately 12, 65 and 90~min); add peak memory alongside runtime in a Supplementary table, which Methods states was recorded}.

\subsection*{A telomere-to-telomere benchmark exposes switch errors hidden by the GRCh38 truth set}
Scoring the same phased VCFs against the two truth sets changed the apparent separation between methods by more than fourfold (Fig.~\ref{fig:benchmarks}). Under GIAB v4.2.1, all methods converged to a narrow band of 1,600--2,200 switch errors genome-wide from 30$\times$ upwards, with corrected \toolname lowest (approximately 1,650), followed by uncorrected \toolname (1,700), \whatshap (1,800), \longphase~1.0 (2,000) and \hapcut (2,100) (Fig.~\ref{fig:benchmarks}b)\todo{exact counts}. Under v5.0q the same VCFs separated four to five times more widely: at 60$\times$ \hapcut made approximately 2,330 switch errors, \whatshap 1,790, \longphase~1.0 1,380, \toolname 743 and \toolname with correction 496, and at 10$\times$ the corresponding counts were 5,030, 4,390, 4,030, 2,258 and 1,517 (Fig.~\ref{fig:benchmarks}c).

Metrics other than the error counts separate the five methods far less. The fraction of benchmark heterozygous SNVs phased (Fig.~\ref{fig:benchmarks}a) and the phase-block N50 (Fig.~\ref{fig:benchmarks}d), which is benchmark independent, differ little between methods\todo{quote the phased fraction and N50 of \longphase~1.0 at 10$\times$ and 60$\times$ (approximately 65\% and 92\%; 0.55 and 2.85~Mb); \longphase~1.0 appears only in Fig.~\ref{fig:benchmarks} and is currently never described quantitatively for these two panels}. Hamming distance was also nearly identical under the two benchmarks for \hapcut, \whatshap and \toolname with and without correction (Fig.~\ref{fig:benchmarks}e,f), indicating that the additional v5.0q switch errors are mostly isolated switches rather than long-range flips. \longphase~1.0 was the exception. Despite an intermediate number of switch errors, \longphase~1.0 showed a Hamming distance of 6--12\% under both benchmarks, which points to long-range flips within blocks. The read-based correction and copy-number-aware filtering introduced in \toolname remove these.

\subsection*{Co-phasing indels, SVs and methylation lengthens blocks by up to 75\% but raises the SNV switch error rate}
Joint phasing of all four variant classes extended haplotypes well beyond the SNV-only limit, at a measurable cost in SNV accuracy that correction then more than recovered (Figs.~\ref{fig:cophasing} and~\ref{fig:combinations}). Indels and methylation sites provide new links across SNV-poor regions, so block N50 at 60$\times$ increased from approximately 2.9~Mb for SNV-only phasing to 4.3~Mb with indels, to 3.1~Mb with methylation alone and to 5.1~Mb with all four classes (Fig.~\ref{fig:combinations}, right column). At the same time, indel observations are noisier than SNV observations: the SNV switch error rate of uncorrected co-phasing rose from 0.034\% (SNV only) to 0.043\% (all classes) at 60$\times$ and from 0.121\% to 0.154\% at 10$\times$, while the Hamming distance rose from 1.8\% to 7.2\% at 60$\times$ (Fig.~\ref{fig:cophasing}e,f). Decomposing the combinations showed that indels carry almost all of the additional error. SNV+SV and SNV+methylation had switch error rates and Hamming distances indistinguishable from SNV-only phasing, while SNV+indel matched the four-class configuration (Fig.~\ref{fig:combinations}).

Adding classes did not make the phasing of any single class less complete. The fraction of phased SNVs was unchanged relative to SNV-only phasing (Fig.~\ref{fig:cophasing}a). The fraction of benchmark heterozygous indels that were phased rose from 31\% at 10$\times$ to 45\% at 60$\times$ and was unaffected by the neural network (Fig.~\ref{fig:cophasing}b)\todo{clarify denominator: all GIAB indels or those called by Clair3}. Of the allele-specific 5mC sites identified by \texttt{modcall}, more than 99\% were phased at every coverage, and the neural network unphased 1.5--5\% of those sites, mostly at low coverage (Fig.~\ref{fig:cophasing}c). Roughly 68--70\% of heterozygous SVs called by Sniffles2 were phased, decreasing slightly with coverage as more low-support SVs were called, and correction removed a further 2 percentage points (Fig.~\ref{fig:cophasing}d).

\todo{Two comparisons are missing from this section and will be requested. (i) \methphaser, which Methods states was run in the methylation-assisted setting, is never reported: add the SNV+5mC comparison of \toolname against \whatshap followed by \methphaser (block N50, switch error rate, Hamming distance) or remove \methphaser from Methods. (ii) \longphase~1.0 already co-phases SNVs with SVs, so the SNV+SV and four-class configurations should be compared against it to separate the contribution of the new evidence classes from that of the new evidence calibration; \longphase~1.0 currently appears only in Fig.~\ref{fig:benchmarks}. Also report the phasing accuracy, not only the phased fraction, of the indels, SVs and 5mC sites themselves (indels and SVs against the v5.0q and GIAB SV benchmarks; 5mC against imprinted differentially methylated regions, for which parent-of-origin is known).}

\subsection*{Neural-network correction reduces switch errors in every variant-class configuration}
The neural network reduced the SNV switch error rate in all five configurations and at every coverage (Fig.~\ref{fig:combinations}, left column). For SNV-only phasing the rate fell by about one-third at every coverage (31--34\%; for example from 0.034\% to 0.023\% at 60$\times$); for all four classes it fell by 37--40\%, from 0.043\% to 0.027\% at 60$\times$ and from 0.154\% to 0.093\% at 10$\times$, and SNV+indel behaved alike. Corrected co-phasing was therefore more accurate than uncorrected SNV-only phasing at every coverage (Fig.~\ref{fig:cophasing}e) while retaining substantially longer blocks (Fig.~\ref{fig:combinations}, right column). The correction was even more effective on the Hamming distance of indel-containing configurations, which fell from 7.2\% to 3.3\% at 60$\times$ and from 11.0\% to 7.0\% at 10$\times$ (Fig.~\ref{fig:combinations}, middle column), which suggests that a large share of the indel-induced errors are long-range flips the network recognizes from the surrounding vote structure. The price was a moderate reduction in block N50 for indel-containing configurations, from 4.3~Mb to 3.7~Mb (SNV+indel) and from 5.1~Mb to 4.1~Mb (all classes) at 60$\times$, because unphased variants that bridged two blocks split those blocks. For configurations without indels the N50 was almost unchanged. On balance, co-phasing remains worthwhile: even after correction, four-class co-phasing yielded blocks approximately 40\% longer than SNV-only phasing at 60$\times$, together with a lower switch error rate than uncorrected SNV-only phasing.

\todo{Ablations a reviewer will demand for this section: (i) the network versus the simple baseline it replaces, namely unphasing every variant with PE above a fixed threshold, matched for the number of variants unphased (this is the control for the claim made in the next section that a fixed entropy threshold cannot substitute for the network); (ii) sensitivity of all metrics to the break threshold (0.30) and to the PE trigger threshold (0.8), ideally as a switch-error-versus-N50 trade-off curve; (iii) precision and recall of the network on held-out switch-error labels. Also state explicitly, here or in Methods, which HG002 chromosomes and coverages were held out of training, since training and evaluation use the same individual.}

\subsection*{Unphased variants are predominantly unsupported heterozygous calls}
To see what the network actually removes, we classified every SNV that \toolname phased and the network subsequently unphased (SNV-only phasing, single replicate, seed 1) against the v5.0q benchmark (Fig.~\ref{fig:unphased}). The network unphased approximately 60,000 SNVs at every coverage, but 86--96\% of those SNVs were calls that the benchmark does not support as heterozygous, either because the site is absent from the benchmark (80--94\%) or because the benchmark calls the site homozygous. Only 4.0\% (2,411 variants) at 60$\times$ and 14\% at 10$\times$ were genuine heterozygous sites (Fig.~\ref{fig:unphased}a). Little real phasing information is therefore lost.

Cross-tabulating each unphased variant by read depth and by phasing informativeness, the number of informative reads linking each variant to its neighbours, showed that genuine heterozygotes were removed overwhelmingly at normal depth with weak phasing evidence (61--81\%), whereas benchmark-absent calls shifted towards low depth as coverage increased (12.5\% at 10$\times$ to 58.6\% at 60$\times$), as expected in hard-to-map regions where reads are lost to alternative loci (Fig.~\ref{fig:unphased}b). High depth alone never caused a variant to be unphased. Among the genuine heterozygotes that were removed, 43\% at 10$\times$ had unanimous but sparse haplotype votes, compared with 17\% of benchmark-absent calls, and this share fell to 22\% at 60$\times$ as evidence accumulated; a further 21--25\% of genuine heterozygotes and 16--18\% of benchmark-absent calls had near-equivocal votes at every coverage (Fig.~\ref{fig:unphased}c)\todo{PE $\ge 0.8$ prevalences read from Fig.~\ref{fig:unphased}c; replace with exact values}. The vote balance at the variant itself would therefore have missed most of the removed heterozygotes; the network, which sees the votes of all neighbours in the window, did not. Clustering within 50~bp and low genotype quality, by contrast, were present in 46--49\% of benchmark-absent calls at every coverage and were far less frequent among genuine heterozygotes, which marks those calls as variant-caller artefacts rather than phasing failures.

Correction therefore loses little true phasing information, about 2,400 SNVs at 60$\times$ out of 2.19 million phased benchmark SNVs\todo{the \texttt{compare} tables put the loss of phased benchmark SNVs at 60$\times$ at 2,055 (2,193,366 before, 2,191,311 after correction), against the 2,411 genuine heterozygotes counted in Fig.~6a; reconcile the two definitions}. Most removals are false heterozygous calls whose spurious phase would otherwise propagate into haplotagged reads and downstream haplotype-aware callers.

\todo{Generalization experiments missing from the Results as a whole, each of which a reviewer is likely to require, and none of which may be asserted without data: (i) a second individual, e.g.\ HG001 or HG005, to show that the network does not encode HG002-specific structure (the model was trained on HG002 and every result above is on HG002); (ii) PacBio HiFi data, since the Discussion claims the graph algorithm supports both platforms and the Introduction contrasts \toolname with HiFi-restricted joint phasing; (iii) a runtime and peak-memory table across coverages, variant-class configurations and thread counts, covering \toolname, \whatshap, \hapcut and \longphase~1.0; (iv) a demonstration that improved phase propagates downstream, e.g.\ haplotagging accuracy or haplotype-aware variant calling, which the Discussion asserts as the motivation for the work.}

\todo{Item (iv) above may already be done. The group's algorithm slides contain a ClairS somatic-calling experiment across six tumour cell lines (COLO829, H1437, H2009, HCC1395, HCC1937, HCC1954) at five simulated purities (20--100\%), in which the germline phaser supplying the haplotagged alignments is varied between \whatshap~2.4, \longphase~1.7.3 and \longphase~2.0 and everything downstream is held fixed. This is exactly the downstream-impact result the Discussion promises, and it is the strongest available answer to ``does a lower switch error rate actually matter?'': on HCC1395 at 100\% purity, somatic SNV F1 rises from 0.8250 with \longphase~1.7.3 to 0.8398 with \longphase~2.0, with the margin widening at low purity. The experiment must be re-run before it can be used: the ``\longphase~2.0'' in those slides is a pre-release build of the 1.7.3 era, not the version submitted here, so its numbers do not describe the method reported in this paper and the improvement may be larger or smaller once repeated. Three further caveats: the comparison against \whatshap was mixed rather than uniformly favourable, so the claim should be framed as ``improves on the \longphase version ClairS currently embeds'' rather than as a win over all phasers; the experiment uses ClairS v0.4.0 in \texttt{ssrs} mode, which must be stated; and this is germline phasing evaluated by its effect on somatic calling, which must be framed so that it is not read as a somatic-phasing claim. The experimental design, however, is sound and directly reusable: hold the somatic caller fixed, vary only the germline phaser supplying the haplotagged alignments, and sweep purity.}

\todo{The same slides contain an HG002 germline comparison whose numbers differ substantially from Fig.~2 (switch error rate nearly flat at 0.056\% from 10$\times$ to 60$\times$, against 0.125\% falling to 0.035\% here). The configurations differ in every respect (DeepVariant v1.9.0 instead of Clair3, the GIAB benchmark instead of v5.0q, \whatshap~2.4 instead of 2.8, and a pre-release build instead of the submitted version), so the two are not comparable and the slide figures must not be mixed into this manuscript. Separately, confirm which \longphase~1 release is the baseline in Fig.~3: the text says 1.0, the slides used 1.7.3, and the two are different programs. Whichever was run, give its exact version and commit in Methods.}
